

\section{Background}

\subsection{ICT in academic education}

Digital technology has changed higher education in stages. Learning management systems (LMS) such as Canvas have centralized course material, communication and submissions. Lecture recording have added a second stage. Recordings of face-to-face lectures are made available online, and students make significant use of them, with some demonstrated benefits for learning outcomes, while lecturers worry about effects on attendance and engagement \parencite{ocallaghan2017}. Recordings also enable the flipped classroom, where students study material before class and apply it during class. A meta-analysis of 114 studies found a small positive effect on learning outcomes but no effect on student satisfaction \parencite{vanalten2019}.

Assessment has also been digitized. In Norway, LMSs such as Canvas and Blackboard are used for ungraded assessment tasks, while dedicated e-exam systems like Inspera Assessment are used for graded ones \parencite{chirumamilla2021}. At NTNU, digital exams grew from 282 submissions in 2016 to 68\% of all examinations in 2024 \parencite{inspera2025}. Core parts of the teaching process are therefore supported by information systems, but the parts that require judgment, such as grading and feedback, remain largely manual.

\subsection{The Generative AI shift}

Earlier systems store and distribute content. Generative AI (GenAI) and large language models (LLMs) can also interpret and produce it, for example by summarizing material, generating feedback and supporting learners. Reviews identify automated feedback and assessment, learning support and the development of critical skills as the main pedagogical uses. They also identify challenges: unreliable or biased output, privacy and academic integrity risks, and overreliance that can weaken students' own critical thinking \parencite{qian2025}. Human oversight is therefore necessary, which motivates the human-in-the-loop design and the ethics and security analysis in later chapters.

In the Gartner Hype Cycle from August 2024 shown in the course lectures, GenAI was at the Peak of Inflated Expectations, with a plateau of productivity forecast in two to five years \parencite{toussaint2025}. This is a reason to be explicit about realistic benefits and limits.

\subsection{An information systems perspective}

An information system is a set of interrelated components (hardware, software, databases, networks, people and procedures) that collect, process, store and disseminate information to support decision making, coordination, control, analysis and visualization. It is a sociotechnical system that must fit the organization's culture and goals \parencite{toussaint2025}. Introducing GenAI at NTNU is therefore an organizational change, not just a tool addition.

The course distinguishes sustaining (incremental) innovation, which gradually improves existing products, services and ways of operating, from disruptive (radical) innovation, which creates new paradigms. It also distinguishes business process re-engineering (radical, top-down redesign) from continuous improvement (gradual, bottom-up) \parencite{toussaint2025}. The proposed solution is mainly process innovation: the sub-processes are redesigned, modelled in BPMN, and supported by GenAI within existing systems (Canvas, Inspera, Studentweb). This is closer to sustaining innovation and continuous improvement than to re-engineering the whole teaching process.

\subsection{Adoption and change}

Success depends on whether people use the system. The Technology Acceptance Model (TAM) explains use through perceived usefulness and perceived ease of use, and UTAUT adds social influence and facilitating conditions. The DeLone and McLean model links information, system and service quality to use, satisfaction and net benefits. Rogers' diffusion of innovation theory describes uneven adoption across innovators, early adopters, majorities and laggards, and Lewin's change model (unfreezing, moving, refreezing) structures the change process \parencite{stair,toussaint2025}. These frameworks are used in the implementation plan to address resistance and literacy risks.
